%% P31 --- Informacja wzajemna
%%
%% Bliźniak: programs/en/P31-mutual-information.tex. Ramka za ramką, makro za
%% makrem: parity.py porównuje oba wydania W KOLEJNOŚCI, więc zdanie zbudowane
%% po polsku musi nieść te same spany matematyczne i te same literały liczbowe
%% w tej samej kolejności. Cyfra zostaje cyfrą, a słowo zostaje słowem.
%%
%% Uwaga tłumaczeniowa: \enquote{sonda} za angielskie \enquote{probe} --- to
%% zwykłe polskie słowo o realnym użyciu w tym sensie, a nie kalka, więc
%% reguła o zostawianiu angielskich terminów ML go nie obejmuje.
%%
%% ODNIESIENIE PRZED MATEMATYKĄ. Angielskie \enquote{Program P24's pair} musi
%% po polsku zaczynać się od odniesienia, bo C4 czyta kolejność.

\program{Informacja wzajemna}
\label{prog:P31}
\index{informacja wzajemna}

\begin{outcomes}
\outcome{1--6}{powiedzieć, czego tabela korelacji parami nie rozstrzyga, i
  nazwać wielkość, która to rozstrzyga}
\outcome{7--10}{powiedzieć, dlaczego ta wielkość nigdy nie jest ujemna, kiedy
  jest zerem i dlaczego jest symetryczna}
\outcome{11--15}{zapisać informację wzajemną jako różnicę entropii, w każdym
  z dwóch kierunków, i powiedzieć, co mierzy każdy z tych dwóch kawałków}
\outcome{16--26}{powiedzieć, w którą stronę myli się estymator wtykowy na
  danych bez zależności i ile pozycji trzeba, żeby przestało to mieć
  znaczenie}
\outcome{27--35}{powiedzieć, co przetwarzanie może, a czego nie może zrobić z
  informacją, i odczytać tę samą nierówność z obu jej końców}
\outcome{36--39}{powiedzieć, co trzeba by zmierzyć, zanim twierdzenie
  interpretowalnościowe o warstwie będzie uzasadnione}
\end{outcomes}

\begin{quiz}
\item Dwie wielkości mają korelację dokładnie zero. Czego się dowiedziałeś o
  tym, czy jedna wyznacza drugą? \teachesat{1--2}
  \answerto{Niczego. Program~\ref{prog:P24} buduje parę o korelacji dokładnie
  zero, w której jedna jest funkcją drugiej; ten program buduje wielkość,
  która to widzi.}
\item Zapisz informację wzajemną jako różnicę entropii. Czy da się to zrobić
  na więcej niż jeden sposób? \teachesat{14--15}
  \answerto{Na dwa: $H(Y) - H(Y \mid X)$ oraz $H(X) - H(X \mid Y)$. To ta
  sama liczba, a kawałki, z których powstaje, nie są tymi samymi
  kawałkami.}
\item Szacujesz informację wzajemną między dwiema wielkościami naprawdę
  niezależnymi, na skończonej próbie. Co dostajesz? \teachesat{16--17}
  \answerto{Liczbę dodatnią, w wartości oczekiwanej, przy każdym rozmiarze
  próby. Prawda wynosi zero, a estymator jej nie zwraca.}
\item Ktoś ulepsza swoją sondę i podaje wyższy wynik. Co to mówi o warstwie?
  \teachesat{29--30}
  \answerto{Nic nowego. Sonda jest przetwarzaniem, więc jej wynik od zawsze
  był ograniczony tym, co warstwa niosła; lepsza sonda to zdanie o sondzie.}
\item A jeśli sonda wypada źle? \teachesat{31--32}
  \answerto{Też nic, w drugą stronę. To ograniczenie jest ograniczeniem od
  dołu i jest luźne o wielkość, której nikt nie mierzy.}
\end{quiz}

% ============================================================
\section{Czego nie mówią dwa brzegi}
\label{sec:P31-margins}
\index{informacja wzajemna!a zerowa korelacja}

\begin{fr}
Program~\ref{prog:P24} zostawił na stole parę i przypiętą do niej obietnicę.

$Z$ przyjmuje $-1$, $0$ lub $1$, każde z prawdopodobieństwem
$\frac{1}{\val{p24.tri.n}}$, a $W = Z^{2}$. Tamten program dowiódł, dokładnie
na ułamkach, że $\Cov(Z, W) = 0$.

Zauważył też, że $W$ jest \emph{funkcją} $Z$, więc zależność jest całkowita.
Podanie $Z$ mówi ci $W$ wprost.

Zanim pójdziesz dalej: do jakiego wniosku o niezależności tej pary uprawnia
cię korelacja równa zero?
\dotline
\nextframe
\end{fr}

\begin{fr}
\ans{Do żadnego}

Korelacja mierzy prostoliniową część zależności, a tu prostoliniowej części
nie ma: dodatnie i ujemne $Z$ ciągną średnią $ZW$ w przeciwne strony i
znoszą się dokładnie. Cała reszta zależności jest dla niej niewidoczna.

To nie jest wada korelacji. To zdanie o tym, co ta liczba w ogóle mierzyła
\dash{} a Program~\ref{prog:P23} powiedział to samo z drugiej strony: tabela
korelacji parami nie rozstrzyga niezależności.

Warto powiedzieć, czym korelacja \emph{nie} jest. Nie jest bezużyteczna i nie
jest zepsuta: na parze, w której zależność naprawdę jest prostą, mierzy
dokładnie to, co trzeba, i mierzy to tanio. Czego nie potrafi, to odpowiedzieć
na pytanie, o które nigdy jej nie proszono, a \enquote{czy te dwie są
niezależne} jest właśnie tym pytaniem.

Potrzebujesz więc wielkości, która odpowiada na pytanie, którym
Program~\ref{prog:P23} \emph{naprawdę} zdefiniował niezależność. Jakie to
było pytanie?
\dotline
\nextframe
\end{fr}

\begin{fr}
\ans{Czy podanie jednej z nich zmienia rozkład drugiej?}

Oto to pytanie zapisane jako arytmetyka.

Jeśli $X$ i $Y$ są niezależne, to rozkład łączny jest iloczynem brzegów,
$p(x, y) = p(x)\,p(y)$, dla każdej pary. Jeśli nie są \dash{} nie jest.

Cała zależność to więc odstęp między dwoma rozkładami na tych samych
zdarzeniach: łącznym, który masz, i iloczynem brzegów, który byś miał.

Cały Program~\ref{prog:P30} spędziłeś na wielkości mierzącej odstęp między
dwoma rozkładami. Co to za wielkość i co powinno wejść tu w jej dwa
argumenty?
\dotline
\nextframe
\end{fr}

\begin{fr}
\ans{Dywergencja, z rozkładem łącznym jako pierwszym i iloczynem brzegów
jako drugim}

\[ I(X; Y) \;=\; \KL{p(x, y)}{p(x)\,p(y)} \]

i to jest definicja. Nazywa się to \textbf{informacją wzajemną} $X$ i $Y$.

Nic nowego nie zostało zdefiniowane. To obiekt z
Programu~\ref{prog:P30} zastosowany do jednej konkretnej pary rozkładów, a
wszystko, co tamten program o nim udowodnił, obowiązuje tu bez słowa więcej.

\begin{note}
Kolejność ma znaczenie, jak zawsze przy tej wielkości, i rozkład łączny idzie
pierwszy. Sekcja~4 mówi o tym, co się dzieje, gdy drugi argument ma lukę,
której pierwszy nie ma, więc warto już teraz wiedzieć, który jest który.
\end{note}
\end{fr}

\begin{fr}
Zastosuj to do pary, którą zostawił Program~\ref{prog:P24}.

Rozkład łączny kładzie $\frac{1}{\val{p24.tri.n}}$ na każdym z $(-1, 1)$,
$(0, 0)$ i $(1, 1)$, a nigdzie indziej nic. Brzegi to $Z$ jednostajne na
trzech wartościach, a $W$ przyjmuje $0$ z prawdopodobieństwem
$\frac{1}{3}$ i $1$ z prawdopodobieństwem $\frac{2}{3}$.

Policz albo przyjmij na jedną ramkę na wiarę: odpowiedź to
$\val{p31.zw.mi}$ nata.

Teraz pytanie, przy którym warto się zatrzymać. Skąd wzięła się ta liczba
\dash{} czyją jest niepewnością?
\dotline
\nextframe
\end{fr}

\begin{fr}
\ans{To dokładnie $H(W)$ \dash{} cała niepewność $W$}

$\val{p31.zw.mi}$ nata, czyli $\val{p31.zw.mi.bits}$ bita, i jest to $H(W)$
co do ostatniego miejsca, które liczy skrypt.

To nie przypadek, a skrypt asertuje tożsamość, nie liczbę: $W$ jest funkcją
$Z$, więc podanie $Z$ nie zostawia \emph{niczego} nieokreślonego w $W$, a
tym, ile dowiadujesz się o $W$ z $Z$, jest zatem wszystko, co $W$ miało.

\begin{trapbox}
\textbf{\enquote{Zerowa korelacja, więc nie ma czego szukać} to teraz dwie
liczby, a nie argument.}

Na tej samej parze, w tej samej ramce: korelacja wynosi $0$, a informacja
wzajemna $\val{p31.zw.mi}$ nata, czyli każdy nat, jaki $W$ posiada.
Program~\ref{prog:P24} powiedział, że macierz korelacji pełna zer nie jest
dowodem, że zbiór cech nie niesie wspólnej informacji. Tyle kosztuje
zignorowanie tamtego zdania.
\end{trapbox}

\mermaidfig{p31-what-a-margin-drops}{Co odrzucają dwa brzegi}{same brzegi}
\end{fr}

% ============================================================
\section{Dlaczego nie może być ujemna}
\label{sec:P31-nonneg}
\index{informacja wzajemna!nieujemność}

\begin{fr}
Wszystko, co udowodnił Program~\ref{prog:P30}, jest teraz twoje za darmo,
bo to jest tamta wielkość, a nie nowa.

Tamten program dowiódł, że dywergencja nigdy nie jest ujemna \dash{}
stosując nierówność Jensena z Programu~\ref{prog:P19} do ilorazu \dash{} i
że jest zerem dokładnie wtedy, gdy oba rozkłady się zgadzają.

Odczytaj oba te fakty z rozkładem łącznym i iloczynem brzegów w dwóch
gniazdach. Co mówią o $I(X; Y)$?
\dotline
\nextframe
\end{fr}

\begin{fr}
\begin{ansblock}
$I(X; Y) \ge 0$ zawsze, a $I(X; Y) = 0$ \textbf{dokładnie wtedy, gdy rozkład
łączny równa się iloczynowi swoich brzegów} \dash{} co jest definicją
niezależności z Programu~\ref{prog:P23}, wypisaną wprost.
\end{ansblock}

Ta wielkość jest więc zerem wtedy i tylko wtedy, gdy oba są niezależne, a
dodatnia w przeciwnym razie. Tej własności korelacja nie miała, a jej
uzyskanie nie wymagało nowej matematyki: to twierdzenie o dywergencji,
odczytane na konkretnej parze argumentów.

Sprawdzone, nie zadeklarowane: na wszystkich $\val{p31.sym.trials}$
niezerowych układach zliczeń tabeli $2 \times 3$ o komórkach ze zbioru
$\{0, 1, 2, 3\}$, unormowanych do rozkładu, ani jeden nie wychodzi ujemny.
\end{fr}

\begin{fr}
Jest jeszcze jedna własność i warto postawić ci ją jako pytanie, zanim
padnie, bo właśnie spędziłeś cały program na wielkości, która jej nie ma.

Dywergencja z Programu~\ref{prog:P30} zdecydowanie \textbf{nie} jest
symetryczna. Cała jej sekcja~3 była o tym, jak bardzo różnią się oba
kierunki, a sekcja~5 o tym, ile kosztuje wybór między nimi.

Czy $I(X; Y)$ to to samo co $I(Y; X)$?
\dotline
\nextframe
\end{fr}

\begin{fr}
\ans{Tak, dokładnie}

I widać to bez liczenia czegokolwiek. Zamiana $X$ i $Y$ nie zamienia dwóch
argumentów niczego: rozkład łączny $p(x, y)$ to ta sama tabela czytana z
drugiej strony, a iloczyn $p(x)\,p(y)$ to ten sam iloczyn. Żaden argument
dywergencji się nie ruszył.

Sprawdzone na tych samych $\val{p31.sym.trials}$ tabelach, z
odpowiedziami zgodnymi co do ostatniego bitu, jaki niesie arytmetyka.

\begin{note}
Warto jasno powiedzieć, czego to \emph{nie} mówi. Dywergencja nadal jest
niesymetryczna; to ta konkretna para argumentów akurat nie zmienia się przy
zamianie. Sekcja~3 pokazuje asymetrię przetrwałą jedno piętro niżej, w
kawałkach, z których tę wielkość da się złożyć.
\end{note}
\end{fr}

% ============================================================
\section{Symetryczna, z niesymetrycznych kawałków}
\label{sec:P31-conditional}
\index{entropia!warunkowa}
\index{informacja wzajemna!jako spadek niepewności}

\begin{fr}
Definicja jest dokładna i nie jest tym, jak ktokolwiek o tej wielkości mówi.
Ludzie mówią \emph{ile znajomość $X$ mówi ci o $Y$}, co brzmi jak zupełnie
co innego.

Łączy jedno z drugim jedna definicja i jest to jedyny nowy obiekt w tym
programie. \textbf{Entropia warunkowa}
\[ H(Y \mid X) \;=\; H(X, Y) \;-\; H(X) \]
to tyle niepewności $Y$, ile zostaje, gdy $X$ jest znane: niepewność pary
pomniejszona o niepewność samego $X$. Pierwszy człon to entropia z
Programu~\ref{prog:P29} z \emph{parą} jako zdarzeniem, co nie wymaga nowej
definicji \dash{} para jest czymś, co zachodzi, a suma z tamtego programu
biegnie po tym, co zachodzi.

Weź znowu parę z Programu~\ref{prog:P24}, z $W = Z^{2}$. Ile wynosi
$H(W \mid Z)$?
\dotline
\nextframe
\end{fr}

\begin{fr}
\ans{Dokładnie $0$}

$W$ jest funkcją $Z$. Gdy $Z$ leży na stole, w $W$ nie zostaje już żadna
niepewność, a arytmetyka zwraca zero, a nie coś małego.

Teraz to samo pytanie w drugą stronę, bo o nim warto pomyśleć:
$H(Z \mid W)$. Podano ci $W$. Ile z $Z$ pozostaje otwarte?
\dotline
\nextframe
\end{fr}

\begin{fr}
\ans{$\val{p31.zw.hzw}$ nata \dash{} zdecydowanie nie zero}

Gdy powiedziano ci, że $W = 1$, wiesz, że $Z$ to $-1$ albo $1$, i nie wiesz
które. Gdy powiedziano ci, że $W = 0$, znasz $Z$ dokładnie. Średnio
$\val{p31.zw.hzw}$ nata z $Z$ przeżywa podanie $W$.

Obie entropie warunkowe są więc tak różne, jak to możliwe: jedna jest
dokładnie zerem, a druga nie.

\begin{trapbox}
\textbf{\enquote{Informacja wzajemna jest symetryczna, więc zależność jest
symetryczna} nie wynika, a ta para jest kontrprzykładem.}

$Z$ wyznacza $W$, a $W$ nie wyznacza $Z$. Ta asymetria jest prawdziwa, mówi
o niej $H(W \mid Z) = 0$ wobec $H(Z \mid W) = \val{p31.zw.hzw}$, a
informacja wzajemna jej nie widzi. Symetria jest własnością wielkości, a nie
zależności, którą ta wielkość mierzy.
\end{trapbox}
\end{fr}

\begin{fr}
Obie te entropie warunkowe prowadzą do tej samej informacji wzajemnej, a to
są dwie linijki arytmetyki, nie nowy fakt. Rozbij logarytm w definicji:
\[ I(X; Y) \;=\; \sum_{x, y} p(x, y) \ln \frac{p(x, y)}{p(x)\,p(y)}
   \;=\; H(X) + H(Y) - H(X, Y) \]
a entropia warunkowa powyżej mówi, że $H(X, Y) - H(X)$ to $H(Y \mid X)$.
Podstaw ją, a potem zrób to samo z literami zamienionymi miejscami:
\[ I(X; Y) \;=\; H(Y) - H(Y \mid X) \;=\; H(X) - H(X \mid Y) \]
\dash{} niepewność, którą miałeś, pomniejszona o tę, która ci zostaje. To
jest zdanie, które ludzie mają na myśli, mówiąc \emph{ile znajomość $X$ mówi
ci o $Y$}, i jest to ta sama liczba co dywergencja.

Sprawdź to na tej parze, od trudniejszej strony. $H(Z)$ wynosi
$\val{p31.zw.hz}$ nata, a $H(Z \mid W)$ \dash{} $\val{p31.zw.hzw}$. Co
zostaje?
\dotline
\nextframe
\end{fr}

\begin{fr}
\ans{$\val{p31.zw.hz} - \val{p31.zw.hzw} = \val{p31.zw.mi}$ nata}

Czyli tyle, ile dała druga strona: $H(W) - 0$, cała niepewność $W$. Jedna
liczba, osiągnięta z dwóch zestawów kawałków, które mają ze sobą niemal nic
wspólnego.

\begin{aibox}
Dlatego ta sama wielkość nosi dwie nazwy w dwóch literaturach i jest tym
samym. \emph{Przyrost informacji} w drzewie decyzyjnym to
$H(\text{etykieta}) - H(\text{etykieta} \mid \text{podział})$, liczony przy
każdym kandydacie na podział i maksymalizowany. To postać prawostronna.
\emph{Informacja wzajemna} w procedurze selekcji cech to postać
dywergencyjna. Biblioteka oferująca obie oferuje jedną liczbę dwa razy, a
obie zgodzą się z dokładnością, jaką ma jej estymator \dash{} co jest
tematem sekcji~4 i jest mniejsze, niż się zakłada.
\end{aibox}

\mermaidfig{p31-two-routes-one-number}{Dwie drogi do tej samej wielkości}{jedna liczba}
\end{fr}

% ============================================================
\section{Co robi z nią skończona próba}
\label{sec:P31-estimation}
\index{informacja wzajemna!obciążenie estymatora}
\index{estymacja!informacji wzajemnej}

\begin{fr}
Wszystko dotąd zakładało, że znasz rozkład łączny. Nie znasz. Masz próbę, a
liczysz informację wzajemną tabeli, którą ta próba tworzy \dash{} estymat
\textbf{wtykowy}.

Program~\ref{prog:P26} spędził cały program na odstępie między estymatorem a
tym, co estymuje, i na dwóch sposobach mylenia się. To jest to samo pytanie
zadane o tę wielkość.

Zacznij od łatwego przypadku i przewidź, zanim pójdziesz dalej. Dwie
wielkości są \emph{naprawdę niezależne}, więc prawda wynosi dokładnie zero.
Bierzesz skończoną próbę i liczysz estymat wtykowy. Co dostajesz?
\dotline
\nextframe
\end{fr}

\begin{fr}
\ans{Liczbę dodatnią, w zasadzie zawsze}

Nie czasem i nie z niewielką szansą pomyłki w drugą stronę: żeby wyszło
zero, tabela próby musi mieć komórki \emph{dokładnie} w proporcji do swoich
własnych brzegów, a skończona próba prawie nigdy tego nie ma.

Wszystko inne to tabela z jakąś widoczną zależnością, a estymat wtykowy tę
zależność raportuje, bo o to został poproszony.

Oto najmniejszy tego przykład. Dwadzieścia rzutów dwiema uczciwymi monetami,
niezależnymi z konstrukcji:

\transcript{p31-plug-in}

Tabela idealnie płaska raportuje zero. Ta przesunięta o jedną sztukę w
każdej komórce raportuje coś innego, i nie jest to mało.
\end{fr}

\begin{fr}
To jedna tabela. Liczy się to, co estymator robi \emph{średnio}, bo tym
właśnie jest obciążenie.

Dla tabeli dwa na dwa o ustalonej sumie możliwych wyników jest skończenie
wiele \dash{} każdy sposób rozłożenia $N$ zliczeń na cztery komórki
\dash{} a każdy ma znane prawdopodobieństwo. Wartość oczekiwaną można więc
policzyć, dodając je wszystkie, zamiast losować próby i uśredniać.

Przy $N = 10$ takich tabel jest $\val{p31.bias.tables10}$; przy $N = 100$
jest ich $\val{p31.bias.tables100}$. Każda jest wyliczona, a wagi sprawdzone
na zgodność z sumą dokładnie $1$.

Zanim padną liczby: to pomiar czy dowód?
\dotline
\nextframe
\end{fr}

\begin{fr}
\ans{Dowód, dla tego $N$}

Nic nie było losowane i nic nie zostało pominięte, więc odpowiedź jest
wartością oczekiwaną, a nie jej oszacowaniem. To rozróżnienie z
Programu~\ref{prog:P14}, a dyscyplina jest ta sama, której użył
Program~\ref{prog:P30} na własnej rodzinie kandydatów: wyliczaj wyczerpująco
tam, gdzie się da, żeby odpowiedź nie była zdaniem o twoim losowaniu.

Oto, co daje to wyliczenie, na danych, których prawdziwa informacja wzajemna
wynosi dokładnie zero.

\begin{center}
\begin{tabular}{rrrr}
\toprule
$N$ & tabel & $\Ex[\text{wtykowy}]$ & wobec $\frac{1}{2N}$ \\
\midrule
$10$ & $\val{p31.bias.tables10}$ & $\val{p31.bias.n10}$ & $\val{p31.bias.ratio10}$ \\
$20$ & $\val{p31.bias.tables20}$ & $\val{p31.bias.n20}$ & $\val{p31.bias.ratio20}$ \\
$50$ & $\val{p31.bias.tables50}$ & $\val{p31.bias.n50}$ & $\val{p31.bias.ratio50}$ \\
$100$ & $\val{p31.bias.tables100}$ & $\val{p31.bias.n100}$ & $\val{p31.bias.ratio100}$ \\
\bottomrule
\end{tabular}
\end{center}

Każda pozycja w trzeciej kolumnie to liczba dodatnia raportowana na danych,
w których nic nie ma.
\end{fr}

\begin{fr}
Ostatnia kolumna to część, przy której warto zwolnić.

Na to obciążenie istnieje standardowa poprawka. Dla tabeli o $a$ wierszach i
$b$ kolumnach odejmuje ona $\frac{(a-1)(b-1)}{2N}$, co przy dwóch na dwa
daje $\frac{1}{2N}$, a kolumna pokazuje prawdziwe obciążenie jako jej
wielokrotność.

Przeczytaj tę kolumnę z góry na dół. Co robi ta poprawka i gdzie robi to
najgorzej?
\dotline
\nextframe
\end{fr}

\begin{fr}
\ans{Zaniża obciążenie, i najbardziej przy najmniejszym $N$}

$\val{p31.bias.ratio10}$ raza poprawka przy $N = 10$, spadając do
$\val{p31.bias.ratio100}$ przy $N = 100$.

Poprawka jest asymptotyczna: wyprowadzono ją dla dużych $N$ i zbiega, co ta
kolumna pokazuje. Konsekwencja jest tą, której nikt nie wypowiada.

\begin{trapbox}
\textbf{\enquote{Zastosowałem standardową poprawkę na obciążenie, więc
liczba jest solidna.}}

Poprawka jest najmniej dokładna dokładnie tam, gdzie obciążenie jest
największe. Przy $N = 100$ zostawia poniżej dwóch procent obciążenia, a
usuwane obciążenie to $\val{p31.bias.n100}$ nata, którego nikt nie
potrzebował usuwać. Przy $N = 10$ zostawia $\val{p31.bias.left10}$ procent
z niego, a to jest ten reżim, dla którego po poprawkę się sięga.

Środek zaradczy, który działa tam, gdzie go nie potrzebujesz, i zawodzi tam,
gdzie potrzebujesz, jest gorszy niż brak środka, bo usuwa powód do
ostrożności.
\end{trapbox}
\end{fr}

\begin{fr}
Teraz przeskaluj to, i przeskaluj uczciwie.

Ta poprawka ma w sobie $(a-1)(b-1)$, więc obciążenie rośnie z liczbą
\emph{komórek} tabeli i maleje z liczbą pozycji. Szesnaście kategorii w
każdą stronę to $\val{p31.big.cells}$ komórek.

Największa informacja wzajemna, jaką mogą mieć dwie szesnastowartościowe
wielkości, to $\ln \val{p31.big.k} = \val{p31.big.max}$ nata.

Oto pytanie, które powinien zadać przegląd projektu. Ile pozycji potrzeba,
zanim obciążenie zejdzie poniżej jednego procenta tego maksimum?
\dotline
\nextframe
\end{fr}

\begin{fr}
\ans{$\val{p31.big.need}$, czyli około $\val{p31.big.per}$ pozycji na
komórkę}

Przy takim rozmiarze próby obciążenie wynosi $\val{p31.big.bias}$ nata,
jeden procent największej wartości, jaką ta wielkość może przyjąć. Poniżej
liczba, którą liczysz, jest po części pomiarem, a po części rozmiarem próby,
w proporcji, której nic na stronie ci nie mówi.

\begin{warning}
\textbf{Ten program podaje tę liczbę w tę stronę, w którą jest ona ważna, a
pierwsza wersja tego nie robiła.}

Oczywisty ruch to puścić wzór przy jakimś małym $N$ \dash{} szesnaście
kategorii w każdą stronę na pięćdziesięciu pozycjach daje przewidywane
$\num{2.25}$ nata obciążenia, czyli większość maksimum, i robi to efektowne
zdanie. Jest też poza własnym reżimem wzoru: pięćdziesiąt pozycji na
$\val{p31.big.cells}$ komórek to jedna piąta pozycji na komórkę, a wzór jest
asymptotyczny w pozycjach wobec komórek. Sonda tego przypadku wykazała błąd
o jedną trzecią.

Efektowne zdanie jest więc niedostępne, a rozmiar próby dostępny, bo
$\val{p31.big.need}$ pozycji na $\val{p31.big.cells}$ komórek mieści się w
reżimie z zapasem. Cytowanie wzoru poza zakresem, dla którego go
wyprowadzono, to ten sam błąd, o którym jest ta sekcja, tylko piętro wyżej.
\end{warning}
\end{fr}

\begin{fr}
\begin{aibox}
Zestaw te dwie liczby z prawdziwym ustawieniem. Sonda podaje informację
wzajemną między aktywacjami warstwy a etykietą lingwistyczną, a zbiór
ewaluacyjny ma kilkaset zdań.

Aktywacje są ciągłe i zostają zbinowane, sklastrowane albo skwantowane
\dash{} co ustala liczbę komórek, a nikt jej nie podaje. Jeśli ta liczba
jest czymś w rodzaju $\val{p31.big.cells}$, a próba to kilkaset, to jesteś
daleko poniżej $\val{p31.big.per}$ pozycji na komórkę, których wzór
wymagał, więc wzór nie stosuje się, a odpowiedź jakościowa jest gorsza, nie
lepsza: jakiekolwiek jest obciążenie, jest rzędu samej wielkości.
Uruchomienie na etykietach bez żadnego związku z wejściem też wyszłoby
dodatnie.

Istnieje sprawdzian, który kosztuje tylko moc obliczeniową, a jest różnicą
między liczbą a wynikiem.
\end{aibox}

Jaki?
\dotline
\nextframe
\end{fr}

\begin{fr}
\ans{Przetasuj etykiety i puść cały potok jeszcze raz}

Etykiety zachowują rozkład, który miały, i tracą każdy związek z wejściem,
więc cokolwiek to uruchomienie zwróci, wyprodukowały twoje biny, twój
rozmiar próby i twój estymator na danych, w których nic nie ma.

Nikt tego prawie nigdy nie robi, a jest tanie.

Teraz pytanie, które ten przepis ukrywa, i warto na nie odpowiedzieć, zanim
pójdziesz dalej. Ile razy je uruchamiasz?
\dotline
\nextframe
\end{fr}

\begin{fr}
\ans{Więcej niż raz, a powyższe wyliczenie mówi dlaczego}

Jedno przetasowane uruchomienie to jedno losowanie z rozkładu, z którego
pochodzą wyliczone tabele. Tabela wypisała jego średnią; to samo wyliczenie
ma resztę tego rozkładu, a on jest szeroki.

Przy $N = 10$ jego rozrzut to $\val{p31.null.sd10}$ nata wobec średniej
$\val{p31.bias.n10}$. Rozrzut jest większy niż wielkość, którą ma mierzyć,
a mniej więcej jedno uruchomienie na dwadzieścia zwraca
$\val{p31.null.q95}$ nata lub więcej \dash{} na danych, w których nic nie
ma. Trzyma się to przy każdym $N$, które obejmuje tabela, nie tylko przy
najmniejszym.

\begin{trapbox}
\textbf{\enquote{Uruchomiłem przetasowaną kontrolę i ją odjąłem.}}

Kontrola jest zmienną losową, a jedno uruchomienie to jedna jej próbka, więc
podawana różnica może wyjść ujemna albo dwa razy większa niż prawdziwa
nadwyżka, a nic w tym uruchomieniu nie powie, co się stało.

Przetasuj zamiast tego wiele razy i odczytaj, gdzie twój zaobserwowany
estymat leży wśród przetasowanych, zamiast odejmować którykolwiek z nich.
Ten ruch pochodzi z Programu~\ref{prog:P27}, pod inną nazwą: zbuduj rozkład
statystyki przy hipotezie zerowej, w której efekt został usunięty, a potem
odczytaj położenie zaobserwowanej wartości w tym rozkładzie. Tam hipotezą
było, że dwa modele są równie dobre, i wyliczono każdy wynik; tutaj jest nią
to, że etykiety nie mają nic wspólnego z wejściem, a uruchomienia są tym
wyliczeniem.
\end{trapbox}

\mermaidfig{p31-bias-is-not-signal}{Skąd bierze się estymat wtykowy}{sygnał czy próba}
\end{fr}

% ============================================================
\section{Przetwarzanie jej nie tworzy}
\label{sec:P31-dpi}
\index{nierówność o przetwarzaniu danych}
\index{informacja wzajemna!nierówność o przetwarzaniu danych}

\begin{fr}
Zostało jedno twierdzenie i to ono wykonuje najwięcej pracy.

Załóżmy, że $Z$ jest liczone z $Y$ i z niczego więcej \dash{} $Y$ to
warstwa, $Z$ to, co raportuje twoja sonda, a sonda widzi warstwę, a nie
wejście. Zapisz to jako $X \to Y \to Z$.

Sonda może być czymkolwiek: klasyfikatorem liniowym, siecią głęboką, tablicą
przeglądową, rzutem monetą ignorującym wejście.

Zgadnij, zanim pójdziesz dalej. Czy $I(X; Z)$ może przekroczyć $I(X; Y)$?
\dotline
\nextframe
\end{fr}

\begin{fr}
\ans{Nie, nigdy, jakakolwiek by ta sonda była}

\[ I(X; Z) \;\le\; I(X; Y) \qquad \text{gdy } X \to Y \to Z \]

To \textbf{nierówność o przetwarzaniu danych}, a nieformalny odczyt jest jej
całością: \emph{przetwarzaniem nie da się stworzyć informacji.} To, co $Z$
wie o $X$, wie przez $Y$, więc nie może wiedzieć więcej niż $Y$.

Sprawdzone tu, a nie założone. Wobec ustalonego rozkładu łącznego niosącego
$\val{p31.dpi.ixy}$ nata wyliczono każdy kanał z $Y$ do $Z$ na wymiernej
siatce \dash{} $\val{p31.dpi.channels}$ z nich \dash{} i ani jeden go nie
przekracza.
\end{fr}

\begin{fr}
To jest twierdzenie o sondach i warto je czytać powoli, bo rozstrzyga
twierdzenie, które ludzie wygłaszają bez przerwy.

Ktoś trenuje lepszą sondę na tej samej warstwie i podaje wyższy wynik.
Warstwa się nie zmieniła i dane się nie zmieniły.

Czego dowiedzieliśmy się o warstwie?
\dotline
\nextframe
\end{fr}

\begin{fr}
\ans{Niczego o warstwie w ogóle}

Nierówność mówi, że wynik każdej sondy przez cały czas był ograniczony przez
$I(X; Y)$. Sonda, która punktuje wyżej, zbliżyła się do sufitu, który był
tam zawsze. To lepsza sonda, a warstwa jest tym, czym była.

\begin{trapbox}
\textbf{\enquote{Nasza mocniejsza sonda odzyskuje więcej składni, więc
reprezentacja koduje składnię bogaciej, niż sądziliśmy.}}

Drugi człon nie wynika z pierwszego. Reprezentacja koduje to, co koduje;
sonda stała się lepsza w odczytywaniu. To są różne zdania o różnych
obiektach, a zmierzone zostało tylko jedno z nich.

Nierówność jest tym, co czyni z tego twierdzenie, a nie przycinkę: żadna
pojemność sondy nie przepchnie wyniku ponad $I(X; Y)$, więc poprawa jest
zawsze ruchem ku sufitowi, a nigdy ruchem sufitu.
\end{trapbox}
\end{fr}

\begin{fr}
Ta sama nierówność ma drugi koniec i to jego nikt nie cytuje.

Wróć do $\val{p31.dpi.channels}$ wyliczonych kanałów \dash{} każdej sondy
w tej rodzinie. Najlepszy z nich to konkretna liczba i nie jest nią
$\val{p31.dpi.ixy}$.

Sięga $\val{p31.dpi.best}$ nata. Wobec warstwy niosącej
$\val{p31.dpi.ixy}$, ile zostawia za sobą najlepsza dostępna sonda w tej
rodzinie?
\dotline
\nextframe
\end{fr}

\begin{fr}
\ans{Więcej niż jedną czwartą}

I to jest ta \emph{najlepsza} z nich. Każda inna sonda w rodzinie zostawia
więcej.

Nierówność czyta się więc dwojako i oba odczyty bywają rutynowo mylone, w
przeciwnych kierunkach, przez tych samych ludzi:

\begin{center}
\begin{tabular}{ll}
\toprule
\textbf{Lepsza sonda punktuje wyżej} & nie mówi nic o warstwie \\
\textbf{Każda sonda punktuje źle} & też nie mówi nic o warstwie \\
\bottomrule
\end{tabular}
\end{center}

Wynik sondy jest \textbf{ograniczeniem od dołu} tego, co niesie warstwa.
Pierwszy wiersz to sufit, do którego się zbliżamy; drugi to luźne
ograniczenie. Nic w wyniku sondowania nie mówi ci, na który z dwóch
patrzysz, a odstęp nie jest podawany, bo nie jest mierzony.
\end{fr}

\begin{fr}
Warto dokładnie powiedzieć, skąd bierze się tu ten odstęp, bo nie jest to
porażka sond.

Wyliczona wyżej rodzina to każdy kanał z trójwartościowego $Y$ na
dwuwartościowe $Z$ na tamtej siatce, a zawiera wszystkie
$\val{p31.dpi.det}$ odwzorowań, które dokonują wyboru wprost, a nie
losowo.

Powiedz, co każde takie odwzorowanie musi zrobić i dlaczego to coś kosztuje.
\dotline
\nextframe
\end{fr}

\begin{fr}
\ans{Musi scalić dwie z trzech wartości, a cokolwiek scali, część tego, co
je rozróżniało, przepada}

Najlepszy wybór traci najmniej; żaden wybór nie traci niczego. Nie jest to
też argument o samej ósemce: najlepszy ze wszystkich
$\val{p31.dpi.channels}$ wyliczonych kanałów okazuje się być jednym z nich,
więc losowanie wyboru nie odkupuje tego, co kosztowało scalenie.

Sonda o wyjściu mniejszym niż wejście jest dokładnie w tej sytuacji, a to
opisuje każdą sondę, jaką ktokolwiek trenuje.
\end{fr}

\begin{fr}
\begin{aibox}
Nierówność rozstrzyga też pytanie o warstwy, i to akurat w stronę, której
ludzie się nie spodziewają.

Sieć liczy każdą warstwę z warstw wcześniejszych i z niczego więcej, więc
$X \to (\text{wszystko przed warstwą } k+1) \to \text{warstwa } k+1$
i nierówność stosuje się \dash{} przez połączenie rezydualne tak samo jak
przez zwykły łańcuch:
\textbf{informacja o wejściu może z głębokością tylko spadać, nigdy rosnąć.}

Brzmi to, jakby przeczyło wszystkiemu, co się mówi o głębszych warstwach
niosących bogatsze cechy, a nie przeczy \dash{} poprawia to, co takie
twierdzenie może znaczyć. Głęboka warstwa nie może wiedzieć więcej o
wejściu. Może być \emph{łatwiejsza do odczytania}, a to jest zdanie o tym, co
potrafi wyciągnąć mała sonda, a nie o tym, co tam jest. To są dwa końce
nierówności z tej sekcji, a zlewanie ich to właśnie ten błąd.
\end{aibox}

\mermaidfig{p31-probe-is-a-floor}{Co ogranicza wynik sondy}{wynik to podłoga}
\end{fr}

% ============================================================
\section{Co by to rozstrzygnęło}
\label{sec:P31-notmeasured}
\index{informacja wzajemna!czego ta książka nie zmierzyła}

\begin{fr}
Ten program był nietypowo surowy wobec pewnej klasy twierdzeń, więc jest ci
winien zdanie o tym, co zrobił, a czego nie.

Zrobił arytmetykę: definicję, dwie tożsamości, dokładną wartość oczekiwaną
po wyliczonych tabelach i nierówność sprawdzoną wobec każdego kanału w
rodzinie. Nic z tego nie angażuje wytrenowanego modelu.

Czego \emph{nie} zrobił, to nie zmierzył żadnej warstwy czegokolwiek.

\begin{warning}
\textbf{Żadna liczba w tym programie nie została policzona z aktywacji
prawdziwego modelu, a żadne z twierdzeń o sondowaniu nie jest pomiarem.}

Arytmetyka jest dokładna i ogranicza to, co pomiar mógłby znaczyć. Nie mówi
ci, co niesie konkretna warstwa, a ta książka nie ma wytrenowanego modelu,
żeby go zapytać. Programy \ref{prog:P08} i \ref{prog:P11} zostawiły ten sam
dług co do rzędu i powiedziały to; to jest w nim trzecia pozycja.
\end{warning}
\end{fr}

\begin{fr}
Oto więc, czego wymagałoby uzasadnienie twierdzenia postaci \emph{warstwa 12
zawiera informację o składni}, a każda pozycja jest czymś, do czego ten
program dał ci powód.

Zatrzymaj się na chwilę i wypisz, o co byś poprosił, zanim pójdziesz dalej.
\dotline
\nextframe
\end{fr}

\begin{fr}
\begin{ansblock}
\begin{enumerate}
\item \textbf{Liczba komórek}, której nikt nie podaje \dash{} binowanie albo
  klastrowanie, które zamieniło ciągłe aktywacje w coś przeliczalnego.
  Sekcja~4: obciążenie rośnie z komórkami i maleje z pozycjami.
\item \textbf{Rozmiar próby wobec tej liczby}, jako pozycje na komórkę.
\item \textbf{Podłoga z przetasowanych etykiet}: ten sam potok puszczony na
  etykietach bez związku z wejściem, wiele razy, i to, gdzie podawany estymat
  leży wśród tych uruchomień. Jedno z nich jest losowaniem z szerokiego
  rozkładu, a nie samą podłogą.
\item \textbf{Jaka była pojemność sondy}, bo sekcja~5 mówi, że wynik jest
  ograniczeniem od dołu, a większa sonda podnosi go, nie ruszając warstwy.
\item \textbf{Co jest z czym porównywane.} Sam wynik jest nieinterpretowalny;
  wynik wobec przetasowanej podłogi albo wobec tej samej sondy na innej
  warstwie \dash{} nie jest.
\end{enumerate}
\end{ansblock}

Żadna z tych pięciu nie jest trudna. Trzecia rozstrzygnęłaby większość
sporów i nie kosztuje nic poza mocą obliczeniową na ponowne puszczenie
potoku.
\end{fr}

\begin{fr}
Ostatnia rzecz, i jest to uczciwa przeciwwaga.

Nic z tego nie mówi, że te twierdzenia są fałszywe. Informacja wzajemna jest
właściwą wielkością do tego pytania, jest jedyną na stole, która widzi to, co
korelacji umyka, a warstwa niosąca informację o składni to zupełnie rozsądna
rzecz dla warstwy.

Program mówi coś węższego i trudniejszego do podważenia: \textbf{zwykły
materiał dowodowy nie odróżnia tego twierdzenia od jego zaprzeczenia.}
Dodatni estymat dostajesz tak czy owak, a uruchomienie, które by je
rozróżniło, jest tym, którego nikt nie robi.

To jest różnica między twierdzeniem fałszywym a twierdzeniem
nieuzasadnionym, a stwierdzane jest tu tylko to drugie.
\end{fr}

\begin{summarybox}
\sumitem{1--6}{\textbf{\result{Zerowa korelacja niczego nie rozstrzyga}}: para z
  Programu~\ref{prog:P24} ma korelację dokładnie $0$ i informację wzajemną
  $\val{p31.zw.mi}$ nata, czyli całą niepewność $W$.}
\sumitem{3--4}{\result{\textbf{Informacja wzajemna jest dywergencją}:
  $I(X; Y) = \KL{p(x,y)}{p(x)p(y)}$}, odstępem między rozkładem łącznym,
  który masz, a iloczynem brzegów, który byś miał przy niezależności.}
\sumitem{7--8}{\textbf{\result{Nigdy więc nie jest ujemna i jest zerem dokładnie przy
  niezależności}}, oba odziedziczone po Programie~\ref{prog:P30} bez nowego
  dowodu.}
\sumitem{9--10}{\textbf{\result{Jest symetryczna}}, bo zamiana obu wielkości zostawia
  oba argumenty dywergencji bez zmian.}
\sumitem{11--15}{\textbf{\result{Kawałki, z których się ją składa, symetryczne nie
  są}}: na tej samej parze $H(W \mid Z) = 0$, a
  $H(Z \mid W) = \val{p31.zw.hzw}$, i obie drogi dochodzą do
  $\val{p31.zw.mi}$.}
\sumitem{16--19}{\textbf{\result{Estymat wtykowy jest dodatni na danych
  niezależnych}}, w wartości oczekiwanej, przy każdym rozmiarze próby
  \dash{} $\val{p31.bias.n10}$ nata przy $N = 10$ z dokładnego wyliczenia
  wszystkich $\val{p31.bias.tables10}$ tabel.}
\sumitem{20--21}{\textbf{\result{Standardowa poprawka zawodzi najbardziej tam, gdzie
  jest najbardziej potrzebna}}: $\val{p31.bias.ratio10}$ raza obok przy
  $N = 10$, zbiegając do $\val{p31.bias.ratio100}$ przy $N = 100$.}
\sumitem{22--23}{\textbf{\result{Szesnaście kategorii w każdą stronę wymaga
  $\val{p31.big.need}$ pozycji}}, żeby obciążenie zeszło poniżej jednego
  procenta maksimum $\val{p31.big.max}$ nata \dash{} około
  $\val{p31.big.per}$ na komórkę.}
\sumitem{24--26}{\textbf{\result{Sprawdzianem jest uruchomienie na
  przetasowanych etykietach, a jedno z nich jest losowaniem, a nie podłogą}}:
  jego rozrzut $\val{p31.null.sd10}$ nata przewyższa własną średnią
  $\val{p31.bias.n10}$ przy $N = 10$, więc tasuj wiele razy i odczytaj, gdzie
  wśród nich leży zaobserwowany estymat.}
\sumitem{27--28}{\result{\textbf{Przetwarzanie nie tworzy informacji}: jeśli
  $X \to Y \to Z$, to $I(X; Z) \le I(X; Y)$}, sprawdzone wobec wszystkich
  $\val{p31.dpi.channels}$ kanałów w rodzinie.}
\sumitem{29--34}{\textbf{\result{Wynik sondy jest więc ograniczeniem od
  dołu}}, a najlepsza sonda w tej rodzinie i tak zostawia za sobą więcej
  niż jedną czwartą. Lepsza sonda nie mówi nic o warstwie, i gorsza też
  nie.}
\sumitem{36--39}{\textbf{Nic tutaj nie zostało zmierzone na prawdziwym
  modelu}, a stawiane twierdzenie mówi, że zwykły materiał dowodowy nie
  odróżnia tamtego twierdzenia od jego zaprzeczenia \dash{} a nie, że jest
  ono fałszywe.}
\end{summarybox}

\canyou

\begin{testexercises}
\item $X$ jest jednostajne na $\{-2, -1, 1, 2\}$, a $Y = \lvert X \rvert$.
  Bez liczenia dywergencji powiedz, ile wynosi $I(X; Y)$ i dlaczego.
  \answerto{$H(Y) = \ln 2$ nata. $Y$ jest funkcją $X$, więc
  $H(Y \mid X) = 0$ i $I(X; Y) = H(Y) - 0$. Ten sam argument co dla pary z
  Programu~\ref{prog:P24}, z $X$ na czterech wartościach zamiast trzech i
  $Y$ na dwóch, tak jak $W$ \dash{} ale równie prawdopodobnych, i dlatego
  $H(Y)$ to $\ln 2$ tam, gdzie $H(W)$ było $\val{p31.zw.mi}$.}
\item Dwie wielkości mają $I(X; Y) = 0$. Ktoś wnioskuje, że ich korelacja
  wynosi zero. Ma rację i czy odwrotność też jest bezpieczna?
  \answerto{Ma rację, a odwrotność nie jest. $I = 0$ to niezależność, która
  wymusza zerową korelację. Zerowa korelacja nie wymusza $I = 0$, i o tym
  jest cała sekcja~1.}
\item Liczysz wtykową informację wzajemną $\num{0.05}$ nata między binarną
  cechą a binarną etykietą na $60$ pozycjach. Co zrobić najpierw?
  \answerto{Porównać ją z obciążeniem na danych niezależnych przy tym
  rozmiarze: tabela z sekcji~4 stawia tam wartość oczekiwaną między
  $\val{p31.bias.n50}$ a $\val{p31.bias.n100}$ nata, więc mniej więcej piąta
  część $\num{0.05}$ jest tym, co estymator zwraca na niczym \dash{} a
  wartość oczekiwana nie jest sufitem, więc ogon tego rozkładu sięga jeszcze
  wyżej. Przetasuj etykiety i puść jeszcze raz wiele razy, a potem podaj,
  gdzie $\num{0.05}$ leży wśród tego, co te uruchomienia zwrócą.}
\item Kolega podaje, że sonda na warstwie 20 punktuje wyżej niż ta sama
  sonda na warstwie 4, i wnioskuje, że warstwa 20 niesie więcej informacji o
  wejściu. Co jest złego w tym wniosku?
  \answerto{Warstwa 20 jest liczona z warstwy 4, więc nierówność o
  przetwarzaniu danych mówi, że nie może nieść więcej informacji o wejściu
  \dash{} tylko mniej albo tyle samo. Zmieniło się to, jak łatwo sonda tego
  rozmiaru potrafi ją wyciągnąć.}
\item Zapisz $I(X; Y)$ na trzy sposoby: jako dywergencję i jako różnicę
  entropii w każdym z dwóch kierunków.
  \answerto{$\KL{p(x,y)}{p(x)p(y)}$, oraz $H(Y) - H(Y \mid X)$, oraz
  $H(X) - H(X \mid Y)$. Wszystkie trzy to ta sama liczba; drugi i trzeci są
  zbudowane z różnych kawałków.}
\item Dlaczego estymatora wtykowego nie da się naprawić samym powiększeniem
  próby, jeśli liczbie komórek wolno rosnąć razem z nią?
  \answerto{Obciążenie idzie jak komórki na pozycje, więc powiększanie obu
  naraz zostawia ten stosunek tam, gdzie był. Spaść muszą pozycje na
  komórkę, i dlatego sekcja~4 stawia ten wymóg właśnie w tę stronę.}
\end{testexercises}


\begin{furtherproblems}
\item Pokaż, że $I(X; X) = H(X)$, i powiedz słowami, co to znaczy.
  \answerto{$H(X \mid X) = 0$, więc $I(X; X) = H(X) - 0 = H(X)$. Wielkość
  mówi ci o sobie wszystko, a \enquote{wszystko} to dokładnie jej własna
  entropia \dash{} dlatego entropię nazywa się czasem samoinformacją.}
\item Program~\ref{prog:P30} pokazał, że jego dywergencja nie jest
  symetryczna. Wyjaśnij w jednym zdaniu, dlaczego ta jest, nie licząc
  niczego.
  \answerto{Zamiana $X$ i $Y$ zostawia oba argumenty dywergencji bez zmian:
  rozkład łączny to ta sama tabela, a iloczyn brzegów ten sam iloczyn.
  Dywergencja nadal jest niesymetryczna; to ta konkretna para argumentów jest
  nieruchoma przy zamianie.}
\item Potok binuje ciągłą aktywację na $\val{p31.big.k}$ kubełków i liczy
  informację wzajemną wobec $\val{p31.big.k}$-wartościowej etykiety. Ktoś
  proponuje podwoić liczbę kubełków, żeby \enquote{rozdzielić więcej
  struktury}. Co to robi z potrzebnym rozmiarem próby?
  \answerto{Rusza się tylko jedna oś: etykieta nadal jest
  $\val{p31.big.k}$-wartościowa, więc komórek robi się z
  $\val{p31.big.cells}$ \dash{} $\val{p31.big.cells.finer}$, a $(a-1)(b-1)$
  z $\val{p31.big.ab}$ \dash{} $\val{p31.big.ab.finer}$. Pozycje potrzebne
  do utrzymania obciążenia na jednym procencie maksimum rosną z
  $\val{p31.big.need}$ do $\val{p31.big.need.finer}$ \dash{} mniej więcej
  dwukrotnie, a nie czterokrotnie. Drobniejsze biny kupują rozdzielczość i
  płacą za nią pozycjami, a zapłata zwykle nie jest wnoszona.}
\item Nierówność o przetwarzaniu danych mówi $I(X; Z) \le I(X; Y)$, gdy
  $X \to Y \to Z$. Podaj $Z$, dla którego zachodzi ona z równością.
  \answerto{Dowolna odwracalna funkcja $Y$ \dash{} choćby przenumerowanie
  wartości. Nic nie zostaje scalone, więc nic nie ginie. To jest też test na
  to, czy krok przetwarzania może cię cokolwiek kosztować: zapytaj, czy da
  się go cofnąć.}
\item Sekcja~4 zmierzyła obciążenie na danych \emph{niezależnych}.
  Uzasadnij, że estymat wtykowy na danych \emph{zależnych} też nie jest
  wiarygodny bez uruchomienia na przetasowanych etykietach.
  \answerto{Obciążenie nie znika, gdy pojawia się prawdziwa zależność; ono
  się do niej dodaje. Dodatni estymat to więc prawdziwa wielkość plus
  nieznana ilość estymatora, a uruchomienie na przetasowanych etykietach jest
  tym, co mierzy to drugie, żeby dało się je odjąć. Bez niego obu nie da się
  rozdzielić.}
\end{furtherproblems}
